\documentclass[10pt]{extarticle}
\usepackage[a4paper,margin=1.5cm]{geometry}
\usepackage[T1]{fontenc}
\usepackage{amsmath,amssymb}
\usepackage{graphicx}
\usepackage[export]{adjustbox}
\usepackage{multicol}
\usepackage{enumitem}
\usepackage{booktabs,array,tabularx}
\setlength{\columnsep}{0.7cm}
\setlength{\parindent}{0pt}
\setlength{\parskip}{3pt}
\pagestyle{empty}
\makeatletter
% figure registry (figdefs.tex is generated by the extract step)
\newcommand{\qfigdef}[4]{\@namedef{qf@file@#1}{#2}\@namedef{qf@w@#1}{#3}\@namedef{qf@wide@#1}{#4}}
\newcommand{\qf@check}[1]{\@ifundefined{qf@file@#1}{\PackageError{score2}{Unknown figure id #1}{}}{}}
\newcommand{\qf@img}[4][]{\includegraphics[width=#2\dimexpr\@nameuse{qf@w@#4}pt\relax,max width=#3,max height=0.45\textheight,#1]{\@nameuse{qf@file@#4}}}
% Full-width questions: finalize wraps every question that contains a wide figure in \qfullbegin...\qfullend,
% so the WHOLE question (stem, figures, options) is set across the page in reading order, then the two columns
% resume. The question's head - its number, stem and first figure (\qfullsplit marks the end of the head) -
% is kept on one page, so a number is never left at a page foot with its figure overleaf; the rest may break
% across pages normally, so no big gaps are left. Without a \qfullsplit, the whole question is kept together
% if it is under 45% of a page. Never write these yourself.
\newif\ifqf@full
\newif\ifqf@boxed
\newsavebox{\qf@box}
\newcommand{\qf@open}{\global\qf@fulltrue\par\global\qf@boxedtrue
  \setbox\qf@box\vbox\bgroup\hsize\textwidth\linewidth\textwidth}
\newcommand{\qf@emit}[1]{\par\egroup\global\qf@boxedfalse
  \ifdim\dimexpr\ht\qf@box+\dp\qf@box\relax<#1\textheight\noindent\usebox{\qf@box}\par\else\unvbox\qf@box\fi}
\newcommand{\qfullsplit}{\ifqf@boxed\qf@emit{0.85}\fi}
\newcommand{\qf@close}{\ifqf@boxed\qf@emit{0.45}\fi\global\qf@fullfalse}
\newcommand{\qfullbegin}{\end{multicols}\qf@open}
\newcommand{\qfullend}{\qf@close\begin{multicols}{2}}
\newcommand{\qfullstart}{\qf@open}
\newcommand{\qfulldone}{\qf@close}
% \qfig[scale]{id}: block figure, centred, directly where it is written (size set by the tool).
% Fallback if a wide figure is not inside a full-width question: break out of the columns in place - no floats.
\newcommand{\qfig}[2][1]{\qf@check{#2}%
  \ifqf@full\qf@blk{#1}{#2}\else
  \ifnum\@nameuse{qf@wide@#2}=1 \expandafter\@firstoftwo\else\expandafter\@secondoftwo\fi
  {\end{multicols}\begin{center}\qf@img{#1}{\textwidth}{#2}\end{center}\begin{multicols}{2}}%
  {\qf@blk{#1}{#2}}\fi}
\newcommand{\qf@blk}[2]{\par\smallskip{\centering\qf@img{#1}{\linewidth}{#2}\par}\smallskip}
% \qfigi[scale]{id}: inline figure (inside option lists / tables / minipages); never breaks columns.
\newcommand{\qfigi}[2][1]{\qf@check{#2}\qf@img[valign=c]{#1}{\linewidth}{#2}}
\makeatother
% structure macros - lint reads these to check completeness and figure placement
\newcommand{\papersec}[1]{\par\bigskip{\centering\Large\bfseries PAPER-#1\par}\medskip}
\newcommand{\partsec}[2]{\par\medskip\noindent{\large\bfseries PART-#1: #2}\par\smallskip}
\newcommand{\sect}[1]{\par\smallskip\noindent\textbf{#1}\par}
\newcommand{\qn}[1]{\par\smallskip\noindent\textbf{#1.}~}
\newcommand{\paraq}[2]{\par\smallskip\noindent\textbf{Paragraph for Questions #1 and #2}\par}
\newenvironment{qoptions}{\begin{itemize}[leftmargin=1.8em,itemsep=1pt,topsep=2pt,label={}]}{\end{itemize}}
\newcommand{\opt}[2]{\item[(#1)] #2}
% figure registry (generated by score2_tools.py finalize)
\qfigdef{p1_phy_s1_q01_1}{Nurture SRG - Test 02 (06 Dec 2023)-img/p1_phy_s1_q01_1.png}{121.3}{0}
\qfigdef{p1_phy_s1_q02_1}{Nurture SRG - Test 02 (06 Dec 2023)-img/p1_phy_s1_q02_1.png}{69.3}{0}
\qfigdef{p1_phy_s1_q04_1}{Nurture SRG - Test 02 (06 Dec 2023)-img/p1_phy_s1_q04_1.png}{182.4}{0}
\qfigdef{p1_phy_s2_q05_1}{Nurture SRG - Test 02 (06 Dec 2023)-img/p1_phy_s2_q05_1.png}{102.3}{0}
\qfigdef{p1_phy_s2_q06_1}{Nurture SRG - Test 02 (06 Dec 2023)-img/p1_phy_s2_q06_1.png}{83.5}{0}
\qfigdef{p1_phy_s2_q07_1}{Nurture SRG - Test 02 (06 Dec 2023)-img/p1_phy_s2_q07_1.png}{105.6}{0}
\qfigdef{p1_phy_s2_q08_1}{Nurture SRG - Test 02 (06 Dec 2023)-img/p1_phy_s2_q08_1.png}{100.4}{0}
\qfigdef{p1_phy_s2_q09_1}{Nurture SRG - Test 02 (06 Dec 2023)-img/p1_phy_s2_q09_1.png}{184.4}{0}
\qfigdef{p1_phy_s2_q10_1}{Nurture SRG - Test 02 (06 Dec 2023)-img/p1_phy_s2_q10_1.png}{90.4}{0}
\qfigdef{p1_phy_s2_q11_1}{Nurture SRG - Test 02 (06 Dec 2023)-img/p1_phy_s2_q11_1.png}{79.1}{0}
\qfigdef{p1_phy_s2_q12_1}{Nurture SRG - Test 02 (06 Dec 2023)-img/p1_phy_s2_q12_1.png}{120.0}{0}
\qfigdef{p1_phy_s3_q02_1}{Nurture SRG - Test 02 (06 Dec 2023)-img/p1_phy_s3_q02_1.png}{103.0}{0}
\qfigdef{p1_phy_s3_q03_1}{Nurture SRG - Test 02 (06 Dec 2023)-img/p1_phy_s3_q03_1.png}{127.5}{0}
\qfigdef{p1_phy_s3_q04_1}{Nurture SRG - Test 02 (06 Dec 2023)-img/p1_phy_s3_q04_1.png}{103.9}{0}
\qfigdef{p1_chem_s1_q04_1}{Nurture SRG - Test 02 (06 Dec 2023)-img/p1_chem_s1_q04_1.png}{93.7}{0}
\qfigdef{p1_chem_s2_q09_A}{Nurture SRG - Test 02 (06 Dec 2023)-img/p1_chem_s2_q09_A.png}{101.1}{0}
\qfigdef{p1_chem_s2_q09_B}{Nurture SRG - Test 02 (06 Dec 2023)-img/p1_chem_s2_q09_B.png}{110.0}{0}
\qfigdef{p1_chem_s2_q09_C}{Nurture SRG - Test 02 (06 Dec 2023)-img/p1_chem_s2_q09_C.png}{102.3}{0}
\qfigdef{p1_chem_s2_q09_D}{Nurture SRG - Test 02 (06 Dec 2023)-img/p1_chem_s2_q09_D.png}{59.2}{0}
\qfigdef{p1_chem_s3_q04_1}{Nurture SRG - Test 02 (06 Dec 2023)-img/p1_chem_s3_q04_1.png}{100.2}{0}
\qfigdef{p1_chem_s3_q06_1}{Nurture SRG - Test 02 (06 Dec 2023)-img/p1_chem_s3_q06_1.png}{59.0}{0}
\raggedcolumns
\makeatletter\@beginparpenalty=10000\makeatother
\newcolumntype{L}[1]{>{\raggedright\arraybackslash}p{#1}}
\begin{document}
\begin{center}{\Large\bfseries Nurture SRG --- Test 02 (06 Dec 2023)}\end{center}
\begin{multicols}{2}
\raggedcolumns
\papersec{1}
\partsec{1}{PHYSICS}
\sect{SECTION-I(i)}
\qn{1} A thin uniform elastic rod of natural length $\ell$, area of cross-section $A$, density $\rho$ and Young's modulus $Y$ is just completely immersed in vertical position in fluid of density $2\rho$ by applying a vertically downward force $(F)$ at its top end. The rod is in equilibrium. Elastic deformation energy stored in the rod is (Neglect lateral strain in the rod and the atmospheric pressure)\par
\qfig{p1_phy_s1_q01_1}
\begin{qoptions}
\opt{A}{$\dfrac{7\rho^2A\ell^3g^2}{6Y}$}
\opt{B}{$\dfrac{9\rho^2A\ell^3g^2}{8Y}$}
\opt{C}{$\dfrac{3\rho^2A\ell^3g^2}{4Y}$}
\opt{D}{$\dfrac{4\rho^2A\ell^3g^2}{3Y}$}
\end{qoptions}
\qn{2} A ball of mass $m$ suspended by massless string of length $\ell$ from fixed point. A particle of mass $2m$ strikes in the direction $\theta=45^\circ$ from horizontal and sticks to it. What should be minimum initial velocity of $2m$ mass, so that the system can complete vertical circle.\par
\qfig{p1_phy_s1_q02_1}
\begin{qoptions}
\opt{A}{$3\sqrt{2g\ell}$}
\opt{B}{$3\sqrt{\dfrac{5g\ell}{2}}$}
\opt{C}{$3\sqrt{5g\ell}$}
\opt{D}{$\sqrt{5g\ell}$}
\end{qoptions}
\qn{3} A metal rod of length $\ell$ \& area of cross-section $A$ is suspended to the celling by one of its end. The density of metal is $\rho$, young modulus of elasticity is $Y$. Choose the correct option
\begin{qoptions}
\opt{A}{The stress at mid point is $\ell\rho g$}
\opt{B}{Extension in rod due to it's own weight $\Delta\ell=\dfrac{\rho g\ell^2}{2Y}$}
\opt{C}{Potential energy stored in rod due to its own weight $U=\dfrac{\rho^2AL^3}{3Y}$}
\opt{D}{The stress at mid point is $2\ell\rho g$}
\end{qoptions}
\qn{4} In the figure shown, mass $M$ is moving with velocity $V$ and all after masses are stationary. All collisions are elastic. To transfer maximum velocity to $m$, all masses should have their values in (All masses are different)\par
\qfig{p1_phy_s1_q04_1}
\begin{qoptions}
\opt{A}{Arithmetic progression}
\opt{B}{Harmonic progression}
\opt{C}{Geometric progression}
\opt{D}{Arithmetic geometric progression}
\end{qoptions}
\sect{SECTION-I(ii)}
\qn{5} Two particles A and B each of mass $m$ are connected by a string of length 1m. The particles are kept on a smooth horizontal plane. The initial separation between the particles is 0.5 m. One of the particles is given a velocity $v_0$ as shown in the figure. Mark the correct options just after string becomes tight\par
\qfig{p1_phy_s2_q05_1}
\begin{qoptions}
\opt{A}{Velocity of particle A is $\dfrac{v_0\sqrt3}{4}$}
\opt{B}{Velocity of particle B is $\dfrac{v_0\sqrt7}{4}$}
\opt{C}{Sum of kinetic energy of A and B is $\dfrac5{16}mv_0^2$}
\opt{D}{Loss in kinetic energy of the system is $\dfrac3{16}mv_0^2$}
\end{qoptions}
\qn{6} Figure shows a solid sphere positioned just below the surface of a fluid and held there. Let the density of the fluid be $\sigma$, force applied by fluid on the upper half of the sphere be $F_1$ and the force on lower half be $F_2$ (excluding the contribution due to atmospheric pressure). Pick correct options(s) :-\par
\qfig{p1_phy_s2_q06_1}
\begin{qoptions}
\opt{A}{$F_1=\dfrac13\pi r^3\sigma g$}
\opt{B}{$F_2=\dfrac53\pi r^3\sigma g$}
\opt{C}{$F_1=\dfrac23\pi r^3\sigma g$}
\opt{D}{$F_2=\dfrac73\pi r^3\sigma g$}
\end{qoptions}
\qn{7} Figure shows a thin tube open at one end rotating with constant angular speed $\omega$. The other end of the tube has a small opening. Tube contains a fluid of density $\rho$ which effuses with a speed $u$ relative to the tube. Let $p_A$ be the pressure at the point in the tube just behind the opening. Assume the atmospheric pressure to be $p_0$. Pick correct option(s) for the instant when $h=\dfrac\ell2$.\par
\qfig{p1_phy_s2_q07_1}
\begin{qoptions}
\opt{A}{$p_A=p_0+\dfrac{\rho\omega^2\ell^2}{4}$}
\opt{B}{$p_A=p_0+\dfrac38\rho\omega^2\ell^2$}
\opt{C}{$u=\dfrac{\omega\ell}2$}
\opt{D}{$u=\dfrac{\sqrt3}2\omega\ell$}
\end{qoptions}
\qn{8} Two identical block A and B, connected by a massless string are placed on a frictionless horizontal plane. A bullet having same mass, moving with speed $u$ strikes block B from behind as shown. If the bullet gets embedded into the block B then choose the correct statements\par
\qfig{p1_phy_s2_q08_1}
\begin{qoptions}
\opt{A}{The velocity of A, B, C after collision is $\dfrac u3$}
\opt{B}{Impulse on A due to tension in the string $\dfrac{2mu}3$}
\opt{C}{Impulse on C due to normal force of collision $\dfrac{-mu}3$}
\opt{D}{Impulse on B due to normal force of collision $\dfrac{2mu}3$}
\end{qoptions}
\qn{9} \qfig{p1_phy_s2_q09_1}
The height of fluid of density $\rho$ is $R$. There is a gate in shape of half circular plate as shown in figure having radius $R$. Find force exerted by the fluid on the gate and its torque about the side PQ. (Neglect atmospheric pressure)
\begin{qoptions}
\opt{A}{$F=2\rho gR^3\left(\dfrac13-\dfrac\pi{16}\right)$}
\opt{B}{$F=2\rho gR^3\left[\dfrac\pi4-\dfrac13\right]$}
\opt{C}{$\tau=2\rho gR^4\left[\dfrac\pi4-\dfrac13\right]$}
\opt{D}{$\tau=2\rho gR^4\left[\dfrac13-\dfrac\pi{16}\right]$}
\end{qoptions}
\qn{10} A uniform metallic ring of mass $m$, radius $R$, cross sectional area `a' and Young's modulus $Y$ is kept on a smooth cone of radius $2R$ and semi vertical angle $45^\circ$ as shown. [Assume that extension in the ring is small]\par
\qfig{p1_phy_s2_q10_1}
\begin{qoptions}
\opt{A}{The tension in the ring is same throughout.}
\opt{B}{The tension in the ring is independent of the radius of the ring.}
\opt{C}{The tension in the ring is $\dfrac{mg}{2\pi}$.}
\opt{D}{Elastic potential energy stored in the ring is $\dfrac{m^2g^2R}{4\pi aY}$.}
\end{qoptions}
\qn{11} A uniform rod of mass $m$ and length $\ell$ is acted upon by two forces each of magnitude $F$ as shown in the figure. The length of the rod is $\ell$, area of cross section is $A$, the young's modulus of the material is $Y$. Select correct alternative(s)\par
\qfig{p1_phy_s2_q11_1}
\begin{qoptions}
\opt{A}{Stress at the mid point of the rod is zero.}
\opt{B}{Stress at $x=\dfrac\ell4$ is $\dfrac F{2A}$.}
\opt{C}{Net change in length of the rod is zero.}
\opt{D}{Total energy stored in the rod is zero}
\end{qoptions}
\qn{12} \qfig{p1_phy_s2_q12_1}
Two blocks of masses $m_1$ \& $m_2$ are attached with a spring of spring constant $K$. $F_1$ and $F_2$ are acting on them as shown in figure. Initially spring is at its natural length. Mark the correct options :-
\begin{qoptions}
\opt{A}{If $\theta=0$ maximum elongation in the spring is $\dfrac{2\left(m_2F_1-m_1F_2\right)}{K\left(m_1+m_2\right)}$}
\opt{B}{If $\theta=0$ maximum elongation in the spring is $\dfrac{2\left(m_2F_1+m_1F_2\right)}{K\left(m_1+m_2\right)}$}
\opt{C}{If $\theta=180^\circ$ maximum elongation in the spring is $\dfrac{2\left(m_2F_1+m_1F_2\right)}{K\left(m_1+m_2\right)}$}
\opt{D}{If $\theta=180^\circ$ maximum elongation in the spring is $\dfrac{2\left(m_2F_1-m_1F_2\right)}{K\left(m_1+m_2\right)}$}
\end{qoptions}
\sect{SECTION-II}
\qn{1} A particle of mass $m_1$ collides elastically with a stationary particle of mass $m_2$. Find the maximum angle through which the striking particle may deviate as a result of the collision. Given $(4m_1=5m_2)$ Write the angle in degrees
\qn{2} The surface mass density (mass per unit area) of a semi circular disc of radius `R' varies as $\sigma=\left|\vec a\times\vec r\right|$, where $\vec a$ is a constant vector along positive $x$ axis and $\vec r$ is position vector with respect to origin. If its centre of mass is at $\left(0,\dfrac{a\pi R}{b}\right)$ then find $(a+b)$. ($a$ and $b$ are coprime)\par
\qfig{p1_phy_s3_q02_1}
\qn{3} In the given tube a liquid of density $\rho$ is filled. If the acceleration of the tube towards right direction such that no liquid will falls out of the tube is $\dfrac{\eta g}5$. Find the value of $\eta$\par
\qfig{p1_phy_s3_q03_1}
\qn{4} A cart with mass $m$ moves horizontally on rails with speed $v$ (fig). The rails go further down and smoothly transition into the new horizontal section located $H$ below. On the lower horizontal section of the rails we have another cart of mass $M$. At what initial trolley speed $v$ (in m/s) will the cart of mass $m$ be able to just rise back to the upper horizontal section again after collision with $M$? There is no friction, and the collision is elastic. ($M=9$ kg, $m=1$ kg, $H=0.8$ m, $g=10$ m/s$^2$)\par
\qfig{p1_phy_s3_q04_1}
\qn{5} A circular ring of radius $R$ \& mass $m$ made by a uniform wire of cross-section area $A$ is rotated about a stationary vertical axis passing through its centre \& perpendicular to the plane of the ring. If young's modulus of material of ring is $Y$ \& increment in radius of ring equals to $\dfrac{mR^2\omega^2}{N\pi AY}$ without breaking find value of $N$. If breaking stress of ring is $\sigma$.
\qn{6} Two wires are made of the same material and have the same volume. However, wire 1 has cross-sectional area $A$ and wire-2 has cross-sectional area $3A$. If the length of wire 1 increases by $\Delta x$ on applying force 1 newton, how much force is needed to stretch wire 2 by the same amount ?
\partsec{2}{CHEMISTRY}
\sect{SECTION-I(i)}
\qn{1} The pH value of 0.1 M solutions of $\mathrm{CH_3COONa}$ (I), $\mathrm{CH_3COOH}$ (II), $\mathrm{CH_3COONH_4}$ (III), NaOH (IV) and HCl (V) is in the order
\begin{qoptions}
\opt{A}{I $<$ II $<$ III $<$ IV $<$ V}
\opt{B}{V $<$ IV $<$ III $<$ II $<$ I}
\opt{C}{V $<$ II $<$ III $<$ I $<$ IV}
\opt{D}{V $<$ II $<$ I $<$ III $<$ IV}
\end{qoptions}
\qn{2} What if the pH of a neutral solution at $37^\circ$C, where $K_w$ equals $2.5\times10^{-14}$ ? ($\log2.5=0.4$)
\begin{qoptions}
\opt{A}{7.0}
\opt{B}{13.6}
\opt{C}{6.8}
\opt{D}{6.6}
\end{qoptions}
\qn{3} The number of hydronium ions in 1 mL of an aqueous solution of pH 12.0 at $25^\circ$C is
\begin{qoptions}
\opt{A}{0.01}
\opt{B}{$10^{-12}$}
\opt{C}{$6.02\times10^8$}
\opt{D}{$6.02\times10^{11}$}
\end{qoptions}
\qn{4} The IUPAC name of\par
\qfig{p1_chem_s1_q04_1}
\begin{qoptions}
\opt{A}{1,1-diethyl-2,2-dimethylpentane}
\opt{B}{4,4-dimethyl-5,5-diethylpentane}
\opt{C}{5,5-diethyl-4, 4-dimethylpentane}
\opt{D}{3-ethyl-4, 4-dimethylheptane}
\end{qoptions}
\sect{SECTION-I(ii)}
\qn{5} If $K_{a_1}$, $K_{a_2}$ and $K_{a_3}$ be the first, second and third dissociation constant of $\mathrm{H_3PO_4}$ and $K_{a_1}\gg K_{a_2}\gg K_{a_3}$ which is/are correct :
\begin{qoptions}
\opt{A}{$\left[\mathrm{H^+}\right]\simeq\sqrt{K_{a_1}\left[\mathrm{H_3PO_4}\right]}$}
\opt{B}{$\left[\mathrm{H^+}\right]\simeq\left[\mathrm{HPO_4^{2-}}\right]$}
\opt{C}{$K_{a_2}\simeq\left[\mathrm{HPO_4^{2-}}\right]$}
\opt{D}{$\left[\mathrm{HPO_4^{2-}}\right]=\left[\mathrm{PO_4^{3-}}\right]$}
\end{qoptions}
\qn{6} Which is/are correct statement(s) about the solubility of AgCl(s) ?\par
Given : $K_{sp(AgCl)}=10^{-10}$; $K_f\left[\mathrm{Ag(NH_3)_2^+}\right]=10^8$.
\begin{qoptions}
\opt{A}{Solubility of AgCl in pure water is $10^{-5}$ mole/litre}
\opt{B}{Solubility of AgCl in 2M KBr is $10^{-5}$ gm/litre}
\opt{C}{Solubility of AgCl in 2 M $\mathrm{AgNO_3}$ is $5\times10^{-11}$ M}
\opt{D}{Solubility of AgCl in 2M $\mathrm{NH_3}$ is 0.166 M}
\end{qoptions}
\qn{7} Degree of hydrolysis (h) for a salt of strong acid and weak base is :
\begin{qoptions}
\opt{A}{Independent of dilution}
\opt{B}{Increases with dilution}
\opt{C}{Increases with decrease in $K_b$}
\opt{D}{Increases with increase in temperature}
\end{qoptions}
\qn{8} $\mathrm{H_2A}$ is a weak diprotic acid. If the pH of 0.1M $\mathrm{H_2A}$ solution is 3 and concentration of $\mathrm{A^{2-}}$ is $10^{-12}$ at $25^\circ$C.\par
Select correct statement(s)
\begin{qoptions}
\opt{A}{$\left[\mathrm{H^+}\right]_{\text{total}}\approx\left[\mathrm{H^+}\right]$ from first step of ionization of acid $\mathrm{H_2A}$}
\opt{B}{Concentration of $\mathrm{OH^-}$ in solution is $10^{-3}$ M}
\opt{C}{The value of $K_{a_1}$ is nearly $10^{-5}$}
\opt{D}{$pK_{a_2}-pK_{a_1}=9$}
\end{qoptions}
\qn{9} Which of the following IUPAC names is/are correct
\begin{qoptions}
\opt{A}{\qfigi{p1_chem_s2_q09_A}\par 3-cyanomethyl-1,5-pentanedinitrile}
\opt{B}{\qfigi{p1_chem_s2_q09_B}\par 3-cyano-6-Hydroxyoct-5-enamide}
\opt{C}{\qfigi{p1_chem_s2_q09_C}\par 5-($2^1$-ethyl butyl)-3,3-dimethyldecane}
\opt{D}{\qfigi{p1_chem_s2_q09_D}\par 2-Bromo-5-hydroxybenzonitrile}
\end{qoptions}
\qn{10} In an acidic indicator HIn has ionization constant is $10^{-8}$. The acid form of indicator is yellow and alkaline form is red. Which is/are correct statements ?\par
(Given : $\log2=0.3$, $\log3=0.48$)
\begin{qoptions}
\opt{A}{The pH range of indicator is 7 to 9}
\opt{B}{Change in pH is 0.96 when 75 \% yellow colour change to 75 \% red colour}
\opt{C}{This indicator is suitable for the titration of strong acid vs. strong base}
\opt{D}{pH of indicator is 8.3 when ratio of acid form to alkaline form is 2.}
\end{qoptions}
\qn{11} A buffer solution can be prepared from a mixture of
\begin{qoptions}
\opt{A}{$\mathrm{NH_4Cl}$ and NaOH in 2:1 mole ratio}
\opt{B}{$\mathrm{CH_3COONa}$ and HCl in 1:1 mole ratio}
\opt{C}{$\mathrm{CH_3COONa}$ and HCl in 2:1 mole ratio}
\opt{D}{$\mathrm{CH_3COONa}$ and HCl in 1:2 mole ratio}
\end{qoptions}
\qn{12} $\mathrm{H_2CO_3}$ ionizes as :\par
$\mathrm{H_2CO_3}\rightleftharpoons\mathrm{H^+}+\mathrm{HCO_3^-}$; $K_1=4.0\times10^{-6}$\par
$\mathrm{HCO_3^-}\rightleftharpoons\mathrm{CO_3^{2-}}$; $K_2=5.0\times10^{-11}$\par
The correct information(s) related with 0.5 M $-$ $\mathrm{Na_2CO_3}$ solution is/are
\begin{qoptions}
\opt{A}{The degree of hydrolysis of $\mathrm{Na_2CO_3}$ is 0.02}
\opt{B}{pH of solution is 2.0}
\opt{C}{pOH of solution is 2.0}
\opt{D}{$\left[\mathrm{H_2CO_3}\right]\approx2.5\times10^{-9}$ M}
\end{qoptions}
\sect{SECTION-II}
\qn{1} The ionic product of heavy water, $\mathrm{D_2O}$, is $10\times10^{-16}$ at $7^\circ$C. The $p^D$ value of pure heavy water at $7^\circ$C is
\qn{2} The self-ionization constant for pure formic acid, $K=\left[\mathrm{HCOOH_2^+}\right]\left[\mathrm{HCOO^-}\right]$ has been estimated as $10^{-6}\,\mathrm{M^2}$ and the density of formic acid is 1.15 g/cm$^3$ at room temperature. If `x $\%$' of formic acid molecules in pure formic acid is converted to formate ions, then the value of `1000x' is
\qn{3} The solubility of $\mathrm{Pb(OH)_2}$ in water is $5\times10^{-6}$M. Calculate the pH of buffer solution in which solubility of $\mathrm{Pb(OH)_2}$ is $5\times10^{-4}$ mol litre$^{-1}$.
\qn{4} 10 mL of $\mathrm{H_2A}$ (weak diprotic acid) solution is titrated against 0.1 M NaOH. $p^H$ of the solution is plotted against volume of strong base added and following observation is made.\par
\qfig{p1_chem_s3_q04_1}
If pH of the solution at first equivalence point is $pH_1$ and at second equivalence point is $pH_2$.\par
Calculate the value of $(pH_2-pH_1)$ at $25^\circ$C\par
Given: For $\mathrm{H_2A}$, $pK_{a_1}=4.6$ and $pK_{a_2}=8$, $\log25=1.4$
\qn{5} The average concentration of $\mathrm{SO_2}$ in the atmosphere over a city on a certain day is 10 ppm, when the average temperature is 298 K. Given that the solubility of $\mathrm{SO_2}$ in water at 298 K is 1.28 moles litre$^{-1}$ and the $pK_a$ of $\mathrm{H_2SO_3}$ is 2.0, estimate the pH of rain on that day.
\qn{6} IUPAC name of following compound is :\par
\textbf{x}-Chloro-\textbf{y}-methoxy cyclohex-\textbf{z}-ene carboxylic acid\par
\qfig{p1_chem_s3_q06_1}
Then, find the value of $\mathbf{[z-(x+y)]}$ is
\partsec{3}{MATHEMATICS}
\sect{SECTION-I(i)}
\qn{1} 4 people sit around a circular table and everyone rolls a dice. The number of ways in which they can roll the dice such that no two people next to each other roll the same number is
\begin{qoptions}
\opt{A}{750}
\opt{B}{630}
\opt{C}{1296}
\opt{D}{None of these}
\end{qoptions}
\qn{2} If the variable takes values $0,1,2,3,\dots,n$ with frequencies ${}^nC_0,{}^nC_1,{}^nC_2,\dots,{}^nC_n$ respectively, then variance is
\begin{qoptions}
\opt{A}{$\dfrac n3$}
\opt{B}{$\dfrac{2n}5$}
\opt{C}{$\dfrac n4$}
\opt{D}{None of these}
\end{qoptions}
\qn{3} Average height and variance of 5 students in a class is 150cm and 18cm$^2$ respectively. A new student whose height is 156 cm is added to the group. The new variance in cm$^2$ is
\begin{qoptions}
\opt{A}{16}
\opt{B}{20}
\opt{C}{22}
\opt{D}{14}
\end{qoptions}
\qn{4} All the five digit numbers in which each successive digit exceeds its predecessor are arranged in ascending order. The 97th number is
\begin{qoptions}
\opt{A}{24589}
\opt{B}{24679}
\opt{C}{24678}
\opt{D}{24689}
\end{qoptions}
\sect{SECTION-I(ii)}
\qn{5} 5 friends can be accommodated in 5 different rooms of a hotel in $N$ ways such that there are not more than 2 friends in any of the rooms, where
\begin{qoptions}
\opt{A}{$N$ is a perfect square}
\opt{B}{Sum of digits of $N$ is 6}
\opt{C}{$N$ is not divisible by 7}
\opt{D}{$N$ has 24 factors}
\end{qoptions}
\qn{6} We say that a triplet of integers $(a,b,c)$ is sorted if $a<b<c$. Suppose there are $N$ number of sorted triplets of positive integers such that $c\le15$ and the greatest common divisor of $a,b,c$ is greater than 1. Then which of the following is INCORRECT?
\begin{qoptions}
\opt{A}{Sum of digits of $N$ is 10}
\opt{B}{Sum of the digits of $N$ is 9}
\opt{C}{Largest single digit factor of $N$ is 6}
\opt{D}{Largest single digit factor of $N$ is 7}
\end{qoptions}
\qn{7} If the median and range of four numbers $\{x,\,y,\,2x+y,\,x-y\}$, where $0<y<x<2y$, are 10 and 28 respectively, then which of the following is/are CORRECT?
\begin{qoptions}
\opt{A}{$x+y=20$}
\opt{B}{$x-y=6$}
\opt{C}{Mean of the numbers is 18}
\opt{D}{Mean of the numbers is 14}
\end{qoptions}
\qn{8} A, B, C are subsets of $\{1,2,3,\dots,9\}$. The number of ordered triplets $(A,B,C)$ such that $A\cap B\cap C=\phi$ is $m$ and number of ordered triplets $(A,B,C)$ such that $A\cap B\cap C=\phi$, $A\cap B\ne\phi$ and $B\cap C\ne\phi$ is $n$, then
\begin{qoptions}
\opt{A}{Number of factors of $m$ is 10}
\opt{B}{$m-n$ is odd number}
\opt{C}{$n$ is a factor of $m$}
\opt{D}{$n$ is a multiple of 6}
\end{qoptions}
\qn{9} 14 identical black balls and 14 identical white balls are to be arranged in a row. The number of ways to arrange them such that on picking any 7 consecutive balls we get atleast one black ball is ${}^{28}c_k-\ell\cdot{}^{21}c_7+{}^{15}c_m$, then
\begin{qoptions}
\opt{A}{$k+\ell$ can be 29}
\opt{B}{$k-m$ can be 12}
\opt{C}{$k=19$}
\opt{D}{$m,k,\ell$ can be three consecutive term of an A.P.}
\end{qoptions}
\qn{10} There are $n$ lines in a plane, no two of which are parallel and no 3 of them are concurrent. Let the plane be divided by $n$ lines in $a_n$ parts, then
\begin{qoptions}
\opt{A}{$a_n=a_{n-1}+(n-1)$}
\opt{B}{$a_n=a_{n-1}+n$}
\opt{C}{$a_6=22$}
\opt{D}{$a_{10}=56$}
\end{qoptions}
\qn{11} The number of selections of 4 letters from the word COLLEGE is equal to
\begin{qoptions}
\opt{A}{24}
\opt{B}{18}
\opt{C}{Coefficient of $x^4$ in $(1+x)^2\left(1+x+x^2\right)^3$}
\opt{D}{Coefficient of $x^4$ in $(1+x)^3\left(1+x+x^2\right)^2$}
\end{qoptions}
\qn{12} Let $N$ be the number of arrangements of all the digits of the number 12345 such that atleast 3 digits will not come in it's original position, then
\begin{qoptions}
\opt{A}{Sum of digits of $N$ is 10}
\opt{B}{Number of ways to resolve $N$ as product of two nos. is 6}
\opt{C}{Harmonic mean of all divisors of $(N+3)$ is 5}
\opt{D}{$N$ is a prime number}
\end{qoptions}
\sect{SECTION-II}
\qn{1} Let $N$ denote the number of 7 digit positive integers having the property that their digits are in non decreasing order. Find the largest single digit factor of $N$
\qn{2} The number of trailing zeros in a number is the number of zeros at the end of the number, for e.g. 100 has 2 trailing zeros and 1001 has none. Find the number of trailing zeros in the constant term of $\left(x+\dfrac1x\right)^{2014}$.
\qn{3} Find the number of 6 - term strictly increasing geometric sequence such that all terms are positive integers and less than 1000
\qn{4} The number of $n$ - letter words consisting only of the letters A, B and C such that A is never immediately followed by B, B is never immediately followed by C and C is never immediately followed by A are 384. Find value of $n$.
\qn{5} The variance about $-2$ and 2 are 18 and 10 respectively. Find standard deviation of the set of observation
\qn{6} Find the number of nine digit numbers divisible by 11 and having sum of digits as 73.
\end{multicols}
\end{document}
